\subsection{Trattazione qualitativa di un particolare potenziale}
\begin{figure}
\centering
\begin{tikzpicture}
  \draw[->] (-5,0) -- (5,0);
  \draw[->] (0,-3) -- (0,1);

  \draw[thick,smooth] plot file {./grafici/potenziale_particolare.dat};
\end{tikzpicture}
\caption{Andamento di un potenziale simmetrico, nullo all'infinito e con un punto di equilibrio stabile.}
\label{particolare}
\end{figure}
Il potenziale (Figura \ref{particolare}) è attrattivo e tale che
\[
V(x) = V(-x), \quad \lim_{x \to \pm \infty} V(x) = 0, \quad V(x) < 0.
\]
Quindi il solito \( \hat{\mathcal{H}} \ket{\psi} = E \ket{\psi} \), cioè 
\[
\psi''(x) = \frac{2m}{\hbar^2} \left( V(x) - E \right) \psi(x).
\]
Da quello che già conosco mi aspetto uno spettro discreto non degenere, cioè degli stati legati (cioè normalizzabili), per quello che ottengo dalla buca infinita di potenziale. Inoltre ci sarà anche una parte continua dello spettro (cioè stati non normalizzabili).

\bigskip
\noindent
\( \boxed{V(x) > E} \) se \( \psi > 0 \), allora \( \psi'' > 0 \); se \( \psi < 0 \), allora \( \psi'' < 0 \) (Figura \ref{potenziale_generico_V_gt_E}).
\begin{figure}
\centering
\begin{tikzpicture}
  \draw (-2,3) -- (-2,-3);
  \node[draw] at (1,3) {\( \psi > 0 \)};
  \draw (-1.8,0.5) -- (1.8,0.5);
  \draw (0,0.2) -- (0,3);
  \draw[thick] (.5,2) .. controls (.75,1.5) and (1,1) .. (1.7,.59);
  \draw[thick] (-1.7,.59) .. controls (-1,1) and (-.75,1.5) .. (-.5,2);

  \draw (-1.8,-2.7) -- (1.8,-2.7);
  \draw (0,-0.2) -- (0,-3);
  \draw[thick] (.5,-2.52) .. controls (.75,-2.5) and (1.25,-2.3) .. (1.7,-1.64);
  \draw[thick] (-.5,-2.52) .. controls (-.75,-2.5) and (-1.25,-2.3) .. (-1.7,-1.64);

  \draw (2,3) -- (2,-3);
  \node[draw] at (5,3) {\( \psi < 0 \)};
  \draw (2.2,2.5) -- (5.8,2.5);
  \draw (4,0.2) -- (4,3);
  \begin{scope}[shift={(4,2.5)}]
    \draw[thick] (.5,-2) .. controls (.75,-1.5) and (1,-1) .. (1.7,-.59);
    \draw[thick] (-.5,-2) .. controls (-.75,-1.5) and (-1,-1) .. (-1.7,-.59);
  \end{scope}

  \draw (2.2,-0.7) -- (5.8,-0.7);
  \draw (4,-0.2) -- (4,-3);
  \begin{scope}[shift={(4,0)}]
    \draw[thick] (.5,-1.08) .. controls (.75,-1.2) and (1.25,-1.5) .. (1.7,-1.96);
    \draw[thick] (-.5,-1.08) .. controls (-.75,-1.2) and (-1.25,-1.5) .. (-1.7,-1.96);
  \end{scope}

  \node[draw] at (-3,3) {\( V > E \)};
  \draw (-5.8,0) -- (-2.2,0);
  \draw (-5.8,-.5) -- (-2.2,-.5);
  \node (E) at (-2.5,-.8) {\(E\)};
  \draw (-4,-1) -- (-4,.5);
  \begin{scope}[scale=0.4, shift={(-10,0)}]
    \draw[thick,smooth] plot file {./grafici/potenziale_particolare.dat};
  \end{scope}
\end{tikzpicture}
\caption{Trattazione qualitativa delle funzioni d'onda per potenziali \( V(x) > E \).}
\label{potenziale_generico_V_gt_E}
\end{figure}
\( \boxed{V(x) < E} \) se \( \psi > 0 \), allora \( \psi'' < 0 \); se \( \psi < 0 \), allora \( \psi'' > 0 \) (Figura \ref{potenziale_generico_V_lt_E}).
\begin{figure}
\centering
\begin{tikzpicture}
  \draw (-2,3) -- (-2,-3);
  \node[draw] at (1,3) {\( \psi > 0 \)};
  \draw (-1.8,0.5) -- (1.8,0.5);
  \draw (0,0.2) -- (0,3);
  \draw[thick] (.5,1.92) .. controls (.75,1.85) and (1.25,1.6) .. (1.7,1.04);
  \draw[thick] (-.5,1.92) .. controls (-.75,1.85) and (-1.25,1.6) .. (-1.7,1.04);

  \draw (-1.8,-2.7) -- (1.8,-2.7);
  \draw (0,-0.2) -- (0,-3);
  \begin{scope}[shift={(4,1)}]
    \draw[thick] (.5,.08) .. controls (.75,.1) and (1.5,.75) .. (1.7,.96);
    \draw[thick] (-.5,.08) .. controls (-.75,.1) and (-1.5,.75) .. (-1.7,.96);
  \end{scope}
  \draw (2,3) -- (2,-3);
  \node[draw] at (5,3) {\( \psi < 0 \)};
  \draw (2.2,2.5) -- (5.8,2.5);
  \draw (4,0.2) -- (4,3);
  \begin{scope}[shift={(0,-0.5)}]
    \draw[thick,smooth] (.5,-2) .. controls (.75,-1.25) and (1.25,-.75) .. (1.7,-.59);
    \draw[thick,smooth] (-.5,-2) .. controls (-.75,-1.25) and (-1.25,-.75) .. (-1.7,-.59);
  \end{scope}

  \draw (2.2,-0.7) -- (5.8,-0.7);
  \draw (4,-0.2) -- (4,-3);
  \begin{scope}[shift={(4,-3)}]
    \draw[thick,smooth] (.5,2) .. controls (.75,1.25) and (1.25,.75) .. (1.7,.59);
    \draw[thick,smooth] (-.5,2) .. controls (-.75,1.25) and (-1.25,.75) .. (-1.7,.59);
  \end{scope}

  \node[draw] at (-3,3) {\( V < E \)};
  \draw (-5.8,0) -- (-2.2,0);
  \draw (-5.8,-.5) -- (-2.2,-.5);
  \node (E) at (-2.5,-.8) {\(E\)};
  \draw (-4,-1) -- (-4,.5);
  \begin{scope}[scale=0.4, shift={(-10,0)}]
    \draw[thick,smooth] plot file {./grafici/potenziale_particolare.dat};
  \end{scope}
\end{tikzpicture}
\caption{Trattazione qualitativa delle funzioni d'onda per potenziali \( V(x) < E \).}
\label{potenziale_generico_V_lt_E}
\end{figure}

\bigskip
\noindent
Se \( V(x) - E > 0 \) sempre, allora ho Figura \ref{V_gt_E_sempre} quindi non c'è alcuna soluzione.
\begin{figure}
\centering
\begin{tikzpicture}
  \draw (-5,0) -- (-1,0);
  \draw (-3,3) -- (-3,-.5);
  \begin{scope}[shift={(-3,1)}]
    \draw[thick] (0,1) .. controls (.5,.5) and (1.5,.25) .. (2,0.13);
    \draw[thick] (0,1) .. controls (-.5,.5) and (-1.5,.25) .. (-2,0.13);
  \end{scope}
  \node (A) at (-3,-1) {non derivabile};

  \draw (1,0) -- (5,0);
  \draw (3,3) -- (3,-.5);
  \begin{scope}[shift={(3,1)}]
    \draw[thick] (-2,1) .. controls (-1,-.5) and (1,-.5) .. (2,1);
  \end{scope}
  \node (B) at (3,-1) {non normalizzabile};
\end{tikzpicture}
\caption{Soluzioni nel caso \(V(x) > E\) sempre. Nessuna di esse è accettabile.}
\label{V_gt_E_sempre}
\end{figure}

\bigskip
\noindent
Se \( V(x) - E > 0 \) o \( V(x) - E < 0 \) a seconda delle regioni, ho Figura \ref{potenziale_variabile}.
\begin{figure}
\centering
\begin{tikzpicture}
  \draw (-5,0) -- (5,0);
  \draw (0,3) -- (0,-3);

  \draw[thick,smooth] plot file {./grafici/potenziale_particolare_1.dat};
  \draw[thick,dashed,smooth] plot file {./grafici/potenziale_particolare_2.dat};
  \draw[thick,dotted,smooth] plot file {./grafici/potenziale_particolare_3.dat};

  \draw[dashed] (-4,-3) -- (-4,3);
  \node (mx0) at (-4.5,-.5) {\(-x_0\)};

  \draw[dashed] (4,-3) -- (4,3);
  \node (x0) at (4.5,-.5) {\(x_0\)};

  \draw[thick, dotted] (.5,-2) -- (2,-2);
  \node (disc) at (1.5,-2.5) {discontinua};
\end{tikzpicture}
\caption{\( V(x) - E > 0 \) o \( V(x) - E < 0 \) a seconda delle regioni.}
\label{potenziale_variabile}
\end{figure}

\bigskip
\noindent
Poiché se \( E \) cresce, cresce anche \( x_0 \), lo stato a minima energia è quello che ha meno oscillazione.

\bigskip
\noindent
Nel caso \( V(x) > E \) sempre, ho soluzioni oscillanti \( \Rightarrow \) spettro continuo, cioè c'è almeno uno stato legato.

\bigskip
\noindent
Quindi ho:
\begin{enumerate}
\item almeno uno stato legato;
\item \( E_0 > V_0 \), quindi \(\Delta x\) è massimo solo localmente e \(\Delta p\) è minimo nello stato fondamentale;
\item spettro non degenere, quindi ogni operatore che commuta con \( \hat{\mathcal{H}} \) è diagonale nella base degli \( \ket{E} \), cioè 
\[
[A, \hat{\mathcal{H}}] = 0 \quad \Rightarrow \quad A \ket{E} = a_E \ket{E}
\]
\end{enumerate}
Infatti se ho \( \psi_1 \) e \( \psi_2 \) tali che
\[
\psi_1'' = \frac{2m}{\hbar^2} \left( V(x) - E \right) \psi_1, \quad \psi_2'' = \frac{2m}{\hbar^2} \left( V(x) - E \right) \psi_2
\]
e sono normalizzabili, ho che
\[
\psi_2 \psi_1'' - \psi_1 \psi_2'' = 0 = \frac{d}{dx} (\psi_2 \psi_1' - \psi_1 \psi_2') \; \Rightarrow \; \psi_2 \psi_1' - \psi_1 \psi_2' = \text{ costante}
\]
e poiché \( \psi_1 \) e \( \psi_2 \) sono normalizzate,
\[
\lim_{x \to \pm \infty} \psi_1(x) = \lim_{x \to \pm \infty} \psi_2(x) = 0
\]
Quindi la costante è 0. Allora 
\[
\frac{\psi_1'}{\psi_1} = \frac{\psi_2'}{\psi_2} \; \Rightarrow \; \frac{d}{dx} \ln \psi_1(x) = \frac{d}{dx} \ln \psi_2(x) \; \Rightarrow \; \ln \psi_1(x) = \ln \psi_2(x) + \text{ cost.} \; \Rightarrow
\]
\[
\; \Rightarrow \psi_1(x) = e^{\text{cost.}} \psi_2(x)
\]
quindi \( \psi_1 \) e \( \psi_2 \) sono la stessa autofunzione e ho spettro non degenere.